\documentclass[11pt]{article}
\usepackage{graphicx}
\usepackage{booktabs}
\usepackage{geometry}
\usepackage{natbib}
\geometry{margin=1in}
\begin{document}
\title{Topology Analysis of Random Peptide Sequence Networks}
\author{}
\date{}
\maketitle
\section*{Methods}
We analyzed monomer and dimer peptide interaction networks using a graphlet-orbit topology assay derived from prior amyloid fibril network classification work. The assay classifies nodes into fibril topology motifs (e.g., ribbon and prism classes) based on 5-node graphlet orbit counts computed with ORCA. For each sequence and state (monomer/dimer), we computed the fraction of nodes assigned to each topology class and used the dominant class (or `none`/`mixed`) as a sequence-level label. We then summarized topology distributions and monomer-to-dimer transitions. \citep{grazioli2019,grazioli2020,hocevar2016}.
\section*{Results}
\begin{figure}[ht]
\centering
\includegraphics[width=0.95\linewidth]{../plots/topology_fraction_boxplots.png}
\caption{Distribution of topology fractions across sequences in monomer and dimer states.}
\end{figure}
\begin{figure}[ht]
\centering
\includegraphics[width=0.85\linewidth]{../plots/dominant_topology_counts.png}
\caption{Dominant topology counts by state.}
\end{figure}
\begin{figure}[ht]
\centering
\includegraphics[width=0.75\linewidth]{../plots/topology_transition_heatmap.png}
\caption{Monomer-to-dimer transitions of dominant topology.}
\end{figure}
\begin{table}[ht]
\centering
\caption{Dominant topology counts by state.}
\label{tab:dominant_counts}
\begin{tabular}{llll}
\hline
state & dominant & count & percent \\
\hline
dimer & none & 483 & 48.3 \\
dimer & double 1,2 2-ribbon & 444 & 44.4 \\
dimer & 1,2 2-ribbon &  43 & 4.3 \\
dimer & mixed &  27 & 2.7 \\
dimer & 1-ribbon &   3 & 0.3 \\
monomer & none & 437 & 43.7 \\
monomer & double 1,2 2-ribbon & 310 & 31.0 \\
monomer & 1,2 2-ribbon & 170 & 17.0 \\
monomer & mixed &  57 & 5.7 \\
monomer & 1-ribbon &  26 & 2.6 \\
\hline
\end{tabular}
\end{table}

\begin{table}[ht]
\centering
\caption{Dominant topology transitions from monomer to dimer.}
\label{tab:transition_counts}
\begin{tabular}{lll}
\hline
monomer & dimer & count \\
\hline
none & none & 293 \\
double 1,2 2-ribbon & double 1,2 2-ribbon & 163 \\
double 1,2 2-ribbon & none & 137 \\
none & double 1,2 2-ribbon & 131 \\
1,2 2-ribbon & double 1,2 2-ribbon & 106 \\
1,2 2-ribbon & none &  37 \\
mixed & double 1,2 2-ribbon &  37 \\
1,2 2-ribbon & 1,2 2-ribbon &  16 \\
mixed & none &  13 \\
1-ribbon & 1,2 2-ribbon &  12 \\
1,2 2-ribbon & mixed &   9 \\
1-ribbon & double 1,2 2-ribbon &   7 \\
none & 1,2 2-ribbon &   7 \\
double 1,2 2-ribbon & 1,2 2-ribbon &   5 \\
double 1,2 2-ribbon & mixed &   5 \\
none & mixed &   5 \\
1-ribbon & mixed &   4 \\
mixed & mixed &   4 \\
1-ribbon & none &   3 \\
mixed & 1,2 2-ribbon &   3 \\
1,2 2-ribbon & 1-ribbon &   2 \\
none & 1-ribbon &   1 \\
\hline
\end{tabular}
\end{table}

\section*{Additional States Metrics}
We computed network-level organization metrics on the same monomer/dimer edgelists as complementary descriptors of self-assembly patterns. These metrics do not define explicit topology classes (e.g., 1-ribbon/2-ribbon); instead, they summarize connectivity structure and its changes between states. Assortativity captures degree-similarity in contacts (more assortative implies clustered, like-degree mixing). The number of communities summarizes interaction compartmentalization. Core statistics (core max-$k$ and core fraction) summarize the strength and size of the mutual-connectivity scaffold. Role entropy and role polarization summarize heterogeneity in graphlet-orbit roles. Operational state flags binarize these metrics to report how frequently networks are assortative, modular, core-periphery-like, or role-polarized.
Across paired sequences, dimerization is associated with a small decrease in assortativity (median $\Delta=-0.018$), a rise in community count (median $\Delta=+1$), a marked reduction in core fraction (median $\Delta=-0.120$; core max-$k$ roughly stable), and a strong shift toward higher role entropy (median $\Delta=+0.307$) with lower role polarization (median $\Delta=-0.071$). Operational flags are largely unchanged for assortative/modular categories, while core-periphery flags increase modestly in the dimer state (mean $\Delta\approx+0.17$).
\begin{figure}[ht]
\centering
\includegraphics[width=0.98\linewidth]{../plots/additional_states/additional_states_metric_distributions_ACS_doublecol.png}
\caption{Additional-states metrics with panel labels (a--n). Left column: raw monomer vs dimer distributions. Right column: per-sequence changes (dimer minus monomer). Panels (a,b) assortativity; (c,d) number of communities; (e,f) core max-$k$; (g,h) core fraction; (i,j) role entropy; (k,l) role polarization; (m,n) core--periphery flag.}
\end{figure}
\bibliographystyle{plainnat}
\bibliography{references}
\end{document}
